% Chapter 7, typeset from lectures/L07: the README is the opener and the closing review, the
% lecture's Appendix A is a section with the same subsections, and its Appendix B, the exercises,
% closes the chapter. The eeprom_driver_demo is typeset straight from its files in §7.1.4.

\chapter{EEPROM}
\label{c7:ch}

{\large\itshape Memory that survives a reset, and costs milliseconds and wear to write.\par}

\begin{chapterbox}{In this chapter}
\begin{itemize}
  \item EEPROM: non-volatile storage that survives a power cycle or reset, and why it's slow and
    wear-limited compared to RAM.
  \item EEPROM registers: \code{EEAR}, \code{EEDR}, \code{EECR}, and the timed write sequence
    (\code{EEMPE} then \code{EEPE}).
  \item \code{eeprom_write_byte()} / \code{eeprom_read_byte()}, and why interrupts are disabled
    around a write.
  \item Why a plain function is the right tool here: there's exactly one EEPROM on the chip.
  \item Worked example: persisting LED1's state across resets, and why the EEPROM write and UART
    report both happen in \code{main()}, never inside an ISR.
\end{itemize}
\end{chapterbox}

This chapter is about the one memory on the chip that keeps its contents when the power goes.
After it you should be able to explain what EEPROM is for and why it's unsuitable for data that
changes every loop iteration, read from and write to EEPROM at a given address, explain why the
EEPROM's write sequence disables interrupts and what could go wrong if it didn't, and explain why
slow, blocking calls like \code{eeprom_write()} and \code{serial_print()} don't belong inside an
ISR.

Read \secref{c7:sec:eeprom}, the substance of the chapter; its worked example extends the serial
demo of \secref{c6:a:6}. Then work the exercises in \secref{c7:sec:exercises}. The Reflection
exercises, and the question the Code exercise asks, are answered in Appendix~\ref{app:answers};
the Code exercise has no printed solution.

% Section 7.1, from lectures/L07/appendix/a_eeprom.md.

\section{EEPROM}
\label{c7:sec:eeprom}\label{c7:app:a}

\subsection{EEPROM: what it is and why}
\label{c7:a:1}
\begin{itemize}
  \item RAM forgets everything the instant power is removed. EEPROM (Electrically Erasable
    Programmable ROM) doesn't: it's a separate, much smaller memory that keeps its contents across
    a reset or power loss, useful for settings, calibration values, or a counter that should survive
    being unplugged.
  \item The ATmega328P has 1024 bytes of EEPROM, addressed \code{0}--\code{1023}, entirely separate
    from flash (where your program lives) and RAM (where your variables live).
  \item Two costs come with that durability: a single byte write takes a few milliseconds
    (eternity, compared to a RAM write), and each EEPROM cell only tolerates roughly 100,000 write
    cycles before it wears out. Don't write to EEPROM in a tight loop or once per iteration of
    \code{main()}; write only when a value actually needs to persist.
\end{itemize}

\subsection{EEPROM registers}
\label{c7:a:2}
Three registers, all specific to the EEPROM peripheral:
\begin{enumerate}
  \item \textbf{\code{EEAR}} (EEPROM Address Register) holds the address, \code{0}--\code{1023}, to
    read from or write to.
  \item \textbf{\code{EEDR}} (EEPROM Data Register) holds the byte being written, or receives the
    byte just read.
  \item \textbf{\code{EECR}} (EEPROM Control Register) controls and reports on the read/write
    process:
  \begin{itemize}
    \item \code{EERE} (EEPROM Read Enable) starts a read.
    \item \code{EEMPE} (EEPROM Master Write Enable) must be set immediately before \code{EEPE},
      it's a safety interlock, not a data bit.
    \item \code{EEPE} (EEPROM Write Enable) starts a write, and stays set until the write
      completes; a new write must wait for it to clear first.
  \end{itemize}
\end{enumerate}
The datasheet requires \code{EEPE} to be set within four CPU cycles of \code{EEMPE} being set, or
the write silently doesn't happen at all. An interrupt firing in that window is a real risk (a few
cycles is nothing), so a write disables interrupts for the handful of instructions between the
two:

\begin{ccode}
while (EECR & (1U << EEPE)) {} // Wait for any earlier write to finish.

EEAR = addr; // Set destination address (0-1023).
EEDR = byte; // Set data to write (one byte).

cli();                 // Disable interrupts during the write sequence.
EECR |= (1U << EEMPE); // EEMPE must be followed by EEPE within 4 cycles.
EECR |= (1U << EEPE);
sei();                 // Re-enable interrupts once the write sequence is complete.
\end{ccode}

\noindent Reading doesn't have this timing constraint, no write is happening, so there's no
equivalent interlock, just the same ``wait for the peripheral to be free'' pattern already familiar
from \code{ADIF} (Chapter~\ref{c4:ch}) and \code{UDRE0} (Chapter~\ref{c6:ch}):

\begin{ccode}
while (EECR & (1U << EEPE)) {} // Wait for any earlier write to finish.
EEAR = addr;                   // Set the address to read from.
EECR |= (1U << EERE);          // Start read operation.
return EEDR;                   // Return the byte at the specified address.
\end{ccode}

\subsection{\code{eeprom_write_byte()} / \code{eeprom_read_byte()}}
\label{c7:a:3}

\begin{ccode}
#include <stdint.h>

#include <avr/interrupt.h>
#include <avr/io.h>

/**
 * @brief Write a byte to EEPROM at the given address.
 *
 * @param[in] addr EEPROM address to write to (0-1023).
 * @param[in] byte Byte to write.
 */
static void eeprom_write_byte(const uint16_t addr, const uint8_t byte)
{
    // Wait for any earlier write to finish.
    while (EECR & (1U << EEPE)) {}

    // Set destination address and data.
    EEAR = addr;
    EEDR = byte;

    // Perform write operation.
    cli();
    EECR |= (1U << EEMPE);
    EECR |= (1U << EEPE);
    sei();
}

/**
 * @brief Read a byte from EEPROM at the given address.
 *
 * @param[in] addr EEPROM address to read from (0-1023).
 *
 * @return Byte stored at that address.
 */
static uint8_t eeprom_read_byte(const uint16_t addr)
{
    // Wait for any earlier write to finish.
    while (EECR & (1U << EEPE)) {}

    // Set address to read from.
    EEAR = addr;

    // Start read operation.
    EECR |= (1U << EERE);

    // Return the byte stored at the given address.
    return EEDR;
}
\end{ccode}

\subsection{Worked example: persisting LED1's state across resets}
\label{c7:a:4}
Extends \secref{c6:a:6}'s button-toggle-report demo: LED1's state now survives a reset or power
cycle, read back from EEPROM in \code{setup()} and written back out every time it changes. The
circuit is the one of Figure~\ref{c6:fig:serial}, unchanged.

\begin{console}
eeprom_driver_demo/
├── Makefile
├── main.c
├── include/
│   └── driver/
│       ├── eeprom.h
│       └── serial.h
└── source/
    └── driver/
        ├── eeprom.c
        └── serial.c
\end{console}

\noindent The driver files (\file{driver/eeprom.h} and \file{driver/eeprom.c}, both printed in full
below) wrap \secref{c7:a:3}'s \code{eeprom_write_byte()}/\code{eeprom_read_byte()} behind a
buffer-and-length interface, \code{eeprom_write()}/\code{eeprom_read()}, the same shape as
Chapter~\ref{c6:ch}'s \code{serial_write()}. \code{main.c} in full follows them.

\subsubsection{The driver in full}
\cfile{lectures/L07/eeprom_driver_demo/include/driver/eeprom.h}
\needspace{8\baselineskip}
\cfile{lectures/L07/eeprom_driver_demo/source/driver/eeprom.c}

\subsubsection{The program}
\cfile{lectures/L07/eeprom_driver_demo/main.c}

\noindent\code{ISR(PCINT0_vect)} only toggles LED1 and sets \code{led1_event}; the EEPROM write and
the UART report both happen in \code{main()}'s loop instead, once per event, never inside the ISR
itself. That's deliberate, not incidental, and it's good practice generally: both
\code{eeprom_write()} (which can block for several milliseconds waiting on \code{EECR}'s
\code{EEPE} bit, \secref{c7:a:2}) and \code{serial_print()} (which blocks byte-by-byte on
\code{UDRE0}) are slow, blocking calls, and an ISR that spends milliseconds inside either one holds
up every other interrupt for that whole time, including the debounce timer this very feature
depends on. Keep ISRs to the minimum work needed (flip a pin, set a flag) and let \code{main()}'s
loop handle anything that can block, the same reasoning \secref{c6:a:6} gave for never calling
\code{serial_print()} from an ISR.

\subsubsection{Building and running}
Build and flash the same way as \secref{c6:a:6}: on Linux, run \code{make} in the demo's
directory, \file{lectures/L07/eeprom_driver_demo} (\code{make flash} to program the board); on
Windows via Microchip Studio, recreate this folder structure inside the project and add
\file{include} as an include directory. The demo's \file{Makefile} is the one printed in
\secref{c6:sec:build}, with the EEPROM driver added to its source files:

\begin{makecode}
SRC_FILES := main.c \
             source/driver/eeprom.c \
             source/driver/serial.c \
\end{makecode}

% Section 7.2, from lectures/L07/README.md: the goals, the evaluation questions, and the next
% lecture.

\needspace{12\baselineskip}
\section{Review}
\label{c7:sec:review}

\needspace{6\baselineskip}
\subsection*{What you should be able to do now}
After this chapter, you should be able to:
\begin{itemize}
  \item Explain what EEPROM is for and why it's unsuitable for data that changes every loop
    iteration.
  \item Read from and write to EEPROM at a given address.
  \item Explain why EEPROM's write sequence disables interrupts, and what could go wrong if it
    didn't.
  \item Explain why slow, blocking calls like \code{eeprom_write()} and \code{serial_print()}
    don't belong inside an ISR.
\end{itemize}

\needspace{6\baselineskip}
\subsection*{Questions to test yourself}
\begin{enumerate}
  \item Why does \code{eeprom_write_byte()} wait on the \code{EEPE} flag before starting a new
    write?
  \item What real AVR datasheet requirement forces \code{EEMPE} to be followed by \code{EEPE}
    within 4 clock cycles, and why does the example disable interrupts for that window?
  \item In \secref{c7:a:4}'s worked example, why does \code{ISR(PCINT0_vect)} only toggle LED1 and
    set a flag, leaving the actual \code{eeprom_write()} and UART report to \code{main()}?
\end{enumerate}

\needspace{6\baselineskip}
\subsection*{Looking ahead}
Chapter~\ref{c8:ch} is about the watchdog timer, another peripheral the chip has exactly one of,
and a hardware safety net against a program that hangs.

% Section 7.3, from lectures/L07/appendix/b_exercises.md.

\section{Exercises}
\label{c7:sec:exercises}\label{c7:app:b}

\begin{aside}
\textbf{How to work these.} The Reflection exercises are answered on paper, from the chapter. The
Code exercise is a program for the board: build and flash it the way \secref{c2:app:c} (Microchip
Studio) or \secref{c2:app:d} (Linux) describes, or run it in the simulator of
Appendix~\ref{sim:ch} if you have no board yet: its simulation harness (\secref{sim:sec:harness})
saves the EEPROM to a file, so what the program stores survives a restart there too.

\textbf{How to check your work.} The Reflection exercises, and the question the Code exercise asks
about its design, are answered in Appendix~\ref{app:answers}. The Code exercise has no printed
solution: the board tells you whether the program does what the exercise describes. Try each
exercise before you read its answer.
\end{aside}

% ------------------------------------------------------------------------------------------
\exerciseset{Concepts}
\label{c7:set:1}

\reflexercise{Waiting before, not after}
\label{c7:ex:1}
{\raggedright
Why does \code{eeprom_write_byte()} wait on the \code{EEPE} flag before it starts, but
\code{eeprom_read_byte()} doesn't need to wait for anything after its read completes?\par}

\reflexercise{The four-cycle window}
\label{c7:ex:2}
\code{EEMPE} must be followed by \code{EEPE} within four CPU cycles or the write silently fails.
Why does that specifically mean interrupts need to be disabled for that window, and what real
failure would you expect if an ISR fired in between and took longer than four cycles to return?

% ------------------------------------------------------------------------------------------
\exerciseset{Program}
\label{c7:set:2}

\codeexercise{Fewer EEPROM writes}
\label{c7:ex:3}
\secref{c7:a:4}'s \file{eeprom_driver_demo} writes to EEPROM on every single button press. Given
\secref{c7:a:1}'s write-cycle warning, is that actually a good idea for a device expected to run for
years? Extend the demo to reduce wear: keep a small, non-persisted counter of presses since the
last EEPROM write, and only call \code{eeprom_write()} once it reaches \code{5}, instead of on every
press. \code{main()}'s event handling already separates ``an event happened'' from ``write it
out'', build on that rather than restructuring it.

